\label{sec:introduction}

Congressional redistricting is, at bottom, an exercise in democratic engineering.
Mapmakers decide---whether deliberately or through algorithmic neutrality---which communities
share a representative, which elections will be competitive, which incumbents face meaningful
challenges, and how far constituents must travel to access their representative's office.
The academic redistricting literature has devoted substantial effort to measuring the partisan
properties of maps (the L-series), evaluating their compactness (the K-series), and analyzing
their VRA compliance (the D-series). What it has devoted less effort to is systematically
linking geometric map properties to the downstream democratic outcomes that redistricting is
constitutionally required to serve.

This gap matters practically. Courts evaluating redistricting plans ask not just whether a
map is partisan but whether it harms democratic representation. State constitutional provisions
guaranteeing ``free and equal elections'' (Pennsylvania), ``fair districts'' (Michigan), or
requiring that districts ``not dilute'' minority voting strength (Voting Rights Act) are
ultimately about outcomes, not geometry. The geometric metrics we can compute from a shapefile
are useful only insofar as they predict the outcomes that the constitutional provisions care about.
If compactness predicts electoral competitiveness and travel distance but not legislative
polarization, then compactness arguments in court are appropriately limited.

The O-series papers (O.1 through O.4) investigate four outcome dimensions, each in dedicated
empirical analysis. This overview paper provides the integrating framework:

\begin{itemize}
\item \textbf{O.1 --- Electoral Competitiveness}: The distribution of competitive seats under
  enacted vs.\ \textsc{bisect} plans, using VEST precinct-level data and a within-state
  difference-in-differences identification across redistricting cycles.
\item \textbf{O.2 --- Voter Turnout}: Precinct-level turnout as a function of district
  competitiveness, instrumented by the redistricting-induced variation in predicted margin.
\item \textbf{O.3 --- Legislative Polarization}: DW-NOMINATE score gaps between the two
  parties' congressional medians, using redistricting-cycle reshuffling as a natural
  experiment.
\item \textbf{O.4 --- Constituent-Representative Distance}: Mean driving time from census
  tract centroid to the representative's district office address, computed via OSRM
  routing API, with Polsby-Popper as the predictor.
\end{itemize}

The organizing prediction across all four papers is stated at the end of this paper
(Section~\ref{sec:prediction}): \textit{compact algorithmic plans predict better outcomes
on all four dimensions, independent of partisan composition}. This prediction is testable,
falsifiable, and has policy implications that extend beyond the typical gerrymandering debate.
Even in a world without partisan gerrymandering, the choice between compact and non-compact
maps would matter for democratic quality.

\subsection{Why Outcomes, Not Just Geometry}

The academic literature on gerrymandering has generated a rich set of geometric measures
(Section~\ref{sec:literature}), but geometric measures are not self-justifying. A court
that invalidates a map because its Polsby-Popper score is 0.12 is making an implicit
empirical claim: that Polsby-Popper scores of 0.12 are associated with democratic harms
that the constitution prohibits. Without the empirical link from geometry to outcomes,
the constitutional argument is incomplete. Courts have sometimes rejected compactness
arguments precisely because the link was asserted but not demonstrated \citep{veith2004}.

The \textsc{bisect} platform provides a natural test bed for establishing this link.
Because \textsc{bisect} maps are generated without partisan data, any differences in
democratic outcomes between \textsc{bisect} plans and enacted plans that survive
controls for partisan geography must be attributable to geometric features of the maps
themselves. This identification logic---partisan geography controlled, geometry varied---is
the backbone of the O-series.

\subsection{The 2020 Redistricting Cycle as Primary Dataset}

The O-series uses the 2020 redistricting cycle as its primary dataset because it is
the first cycle for which all of the following are simultaneously available:
(1) \textsc{bisect} algorithmic plans for all 50 states; (2) VEST 2020 precinct-level
election returns; (3) ACS 2019--2021 5-year estimates for tract-level demographic covariates;
(4) LODES commuting data for travel distance validation; and (5) enacted plans that can
be compared to algorithmic alternatives using a controlled before-after design.

The 2010 cycle serves as the pre-period for difference-in-differences estimates in O.1 and O.2.
The 2000 cycle is used for robustness checks where VEST data coverage allows.
